\chapter{Einleitung}
\label{cha:Einleitung}

Die Qualität der Raumluft hat einen erheblichen Einfluss auf Gesundheit, Wohlbefinden und Leistungsfähigkeit. Erhöhte CO\textsubscript{2}-Konzentrationen\index{CO2} führen zu Müdigkeit und Konzentrationsverlust, Feinstaub\index{Feinstaub} belastet die Atemwege und flüchtige organische Verbindungen (\acrshort{acr:VOC})\index{VOC} können Kopfschmerzen und Schleimhautreizungen verursachen. Obwohl diese Zusammenhänge wissenschaftlich gut belegt sind, verfügen die wenigsten Innenräume über eine kontinuierliche Überwachung der Luftqualität \autocite{WHO.2021}.

Kommerzielle Luftqualitätssensoren sind entweder auf einzelne Parameter beschränkt oder als Komplettsysteme kostspielig. Im Rahmen der Projektarbeit im Modul Mikrocomputertechnik~1 an der \acrshort{acr:DHBW} Ravensburg unter Betreuung von Thomas Stingl (Airbus) wurde daher das System \textbf{InspectAir} entwickelt -- ein kostengünstiges, ESP32-S3-basiertes Messgerät, das fünf relevante Luftqualitätsparameter gleichzeitig erfasst und auf einem 4-Zoll-Display mit intuitiver Farbcodierung darstellt.

\textbf{Ziel} der Arbeit ist die vollständige Entwicklung eines funktionsfähigen Prototyps nach dem V-Modell\index{V-Modell}: von der Anforderungsanalyse über den Systementwurf und die Implementierung bis zur Verifikation durch definierte Testfälle. Das Projekt deckt dabei Hardware-Design, Firmware-Entwicklung und mechanische Konstruktion ab.

Die \textbf{Aufgabenstellung}\index{Aufgabenstellung} fordert die Darstellung eines nahezu vollständigen Entwicklungszyklus mit den Artefakten Requirements Specification, Design Description, Coding Rules, Configuration Management, Schedule, Risk \& Opportunities sowie einer Testspezifikation mit Testergebnissen.

\textbf{Gliederung:} Kapitel~\ref{cha:Grundlagen} führt die technischen Grundlagen ein. Kapitel~\ref{cha:vorgehen} beschreibt das Vorgehen mit Anforderungsanalyse und Systementwurf. Kapitel~\ref{cha:umsetzung} dokumentiert die Hardware- und Software-Umsetzung sowie die Testergebnisse. Kapitel~\ref{cha:zusammenfassung} fasst die Ergebnisse zusammen und gibt einen Ausblick.
